\documentclass[11pt,a4paper]{article}

% --- Encoding ---
\usepackage[utf8]{inputenc}
\usepackage[T1]{fontenc}
\usepackage[spanish,es-noshorthands]{babel}
\usepackage{lmodern}

% --- Layout ---
\usepackage[margin=2.5cm]{geometry}
\usepackage{parskip}
\setlength{\parindent}{0pt}
\setlength{\parskip}{0.5em plus 0.2em minus 0.1em}

% --- Tables and graphics ---
\usepackage{booktabs}
\usepackage{array}
\usepackage{enumitem}
\usepackage{xcolor}
\usepackage{graphicx}
\usepackage{tikz}
\usetikzlibrary{positioning,shapes.geometric,arrows.meta,fit,backgrounds,calc,decorations.pathreplacing}
\usepackage{caption}
\usepackage{amsmath,amssymb}

% --- Code listings ---
\usepackage{listings}
\definecolor{codegray}{rgb}{0.96,0.96,0.96}
\definecolor{codeborder}{rgb}{0.85,0.85,0.85}
\definecolor{codecomment}{rgb}{0.30,0.55,0.30}
\definecolor{codekeyword}{rgb}{0.10,0.30,0.70}
\definecolor{codestring}{rgb}{0.65,0.15,0.10}

\lstdefinestyle{cstyle}{
    language=C,
    backgroundcolor=\color{codegray},
    basicstyle=\ttfamily\footnotesize,
    breaklines=true,
    frame=single,
    rulecolor=\color{codeborder},
    showstringspaces=false,
    keywordstyle=\color{codekeyword}\bfseries,
    commentstyle=\color{codecomment}\itshape,
    stringstyle=\color{codestring},
    aboveskip=1em,
    belowskip=1em,
}

\lstdefinestyle{textstyle}{
    backgroundcolor=\color{codegray},
    basicstyle=\ttfamily\footnotesize,
    breaklines=true,
    frame=single,
    rulecolor=\color{codeborder},
    showstringspaces=false,
    aboveskip=1em,
    belowskip=1em,
    columns=fullflexible,
}

\lstset{style=textstyle}

% --- Definition / Idea / Important boxes ---
\usepackage[most]{tcolorbox}
\newtcolorbox{definicion}[1][]{
    colback=blue!5!white,
    colframe=blue!50!black,
    fonttitle=\bfseries,
    title=Definición,
    sharp corners=southwest,
    rounded corners=northwest,
    arc=2pt,
    boxrule=0.6pt,
    left=8pt,right=8pt,top=4pt,bottom=4pt,
    #1
}
\newtcolorbox{idea}[1][]{
    colback=orange!5!white,
    colframe=orange!60!black,
    fonttitle=\bfseries,
    title=Idea intuitiva,
    sharp corners=southwest,
    rounded corners=northwest,
    arc=2pt,
    boxrule=0.6pt,
    left=8pt,right=8pt,top=4pt,bottom=4pt,
    #1
}
\newtcolorbox{importante}[1][]{
    colback=red!5!white,
    colframe=red!50!black,
    fonttitle=\bfseries,
    title=Importante,
    sharp corners=southwest,
    rounded corners=northwest,
    arc=2pt,
    boxrule=0.6pt,
    left=8pt,right=8pt,top=4pt,bottom=4pt,
    #1
}

% --- Hyperlinks ---
\usepackage[colorlinks=true,
            linkcolor=blue!60!black,
            urlcolor=blue!50!black,
            citecolor=blue!60!black]{hyperref}

% --- Memory diagram color palette ---
\definecolor{mem1}{HTML}{4C72B0}
\definecolor{mem2}{HTML}{DD8452}
\definecolor{mem3}{HTML}{55A868}
\definecolor{mem4}{HTML}{8172B2}
\definecolor{memshared}{HTML}{4FC3C7}
\definecolor{memfile}{HTML}{F4A261}

% --- Title ---
\title{
    \vspace{-2.5em}
    \textbf{Resumen teórico --- Sistemas Operativos} \\[0.3em]
    \large Notas del curso: \emph{Memoria virtual}
}
\author{
    Tomás Rodeghiero \\
    \small Sistemas Operativos --- Universidad Nacional de Río Cuarto
}
\date{2026}

\begin{document}

\maketitle

\begin{abstract}
\noindent
Este documento es un resumen teórico del capítulo \emph{Memoria virtual}
de las \emph{Notas del curso} de Sistemas Operativos (Marcelo Arroyo,
UNRC). Se desarrollan las técnicas que extienden la memoria física
disponible usando disco: \emph{swapping} o \emph{demand paging},
mapeo de archivos en memoria (\texttt{mmap}), carga de programas bajo
demanda, crecimiento dinámico del stack, memoria compartida entre
procesos, \emph{copy on write} en \texttt{fork()} y los algoritmos
clásicos de reemplazo de páginas (FIFO, second chance / clock, LRU /
NFU, NRU, working set). Cierra con el caso de estudio del kernel de
GNU/Linux. El objetivo es servir de material de estudio para preparar
los prácticos asociados a memoria virtual y el segundo parcial.
\end{abstract}

\tableofcontents
\bigskip

% =====================================================================
\section{Introducción: memoria virtual y swapping}
% =====================================================================

La evolución del hardware permite que las CPUs modernas soporten
\emph{enormes} espacios de direcciones: una CPU moderna puede manejar
direcciones de 64\,bits (128\,terabytes) o más. Esto habilita el
desarrollo de programas grandes.

\begin{importante}
En un sistema moderno, la suma de memoria RAM requerida por los procesos
suele ser \textbf{mayor} a la disponible físicamente en el hardware. La
\textbf{memoria virtual} es la técnica que permite ejecutar esos
programas, manteniendo en RAM sólo las partes en uso en cada momento.
\end{importante}

\paragraph{Idea central.}
Un proceso no necesita estar cargado en memoria \emph{completamente}:
sólo necesita las partes que está usando ahora. Con el esquema de
paginado visto en el capítulo anterior, es posible cargar
(\emph{swap-in}) o desalojar (\emph{swap-out}) \emph{partes} de un
proceso con la granularidad del tamaño de página.

\subsection{Swap-in y swap-out}

\begin{itemize}[leftmargin=*]
    \item Un \textbf{swap-in} se produce cuando el proceso referencia
        una página que reside en disco: hay que traerla a RAM.
    \item Un \textbf{swap-out} se produce cuando el sistema necesita
        liberar una página de RAM (por ejemplo, porque no quedan
        páginas libres). Comúnmente lo invoca el \emph{memory
        allocator}.
\end{itemize}

Este proceso de cargar/desalojar páginas en disco se conoce como
\textbf{swapping} o \textbf{demand paging}.

\subsection{Estado de cada página}

Para implementar memoria virtual, por cada página se debe mantener su
\emph{estado}: si está mapeada a una página física en RAM, o si está
guardada en un bloque de disco. Esto se hace almacenando en la
\emph{pte}:

\begin{itemize}[leftmargin=*,nosep]
    \item el \textbf{physical page number} (\texttt{ppn}) si la página
        está en RAM, con el bit \textbf{V} (\emph{valid} o
        \emph{present} en algunas arquitecturas) en 1, o bien
    \item el \textbf{disk block number} (\texttt{dbn}) si la página
        reside en disco, con el bit V en 0.
\end{itemize}

Con esta convención, el hardware genera automáticamente un
\emph{page fault} cuando se referencia una página con $V=0$ y le da al
kernel la oportunidad de hacer el swap-in correspondiente.

\subsection{Extender la RAM con disco}

\begin{center}
\begin{tikzpicture}[
    every node/.style={font=\ttfamily\small},
    cell/.style={draw,minimum height=0.7cm,minimum width=3.0cm,align=center},
    pcell/.style={draw,minimum height=0.7cm,minimum width=2.2cm,align=center}
]
% vpn / page table
\node[font=\bfseries\small] at (0,5.4) {vpn};
\node[font=\bfseries\small] at (2.4,5.4) {page table};
\node[anchor=east] at (-1.1,4.6) {0};
\node[cell] (e0) at (2.4,4.6) {\ldots};
\node[anchor=east] at (-1.1,3.9) {i};
\node[cell,fill=mem3!20] (ei) at (2.4,3.9) {ppn=1\hspace{4pt}V=1};
\node[cell] (e2) at (2.4,3.2) {\ldots};
\node[anchor=east] at (-1.1,2.5) {j};
\node[cell,fill=memfile!30] (ej) at (2.4,2.5) {dbn=k\hspace{4pt}V=0};
\node[cell] (e3) at (2.4,1.8) {\ldots};
\node[cell] (e4) at (2.4,1.1) {\ldots};
\node[anchor=east] at (-1.1,0.4) {n-1};
\node[cell] (en) at (2.4,0.4) {\ldots};

% physical pages
\node[font=\bfseries\small] at (7.5,5.4) {Physical pages};
\node[anchor=east] at (6.0,4.6) {0};
\node[pcell,fill=memshared!20] (f0) at (7.5,4.6) {};
\node[anchor=east] at (6.0,3.9) {1};
\node[pcell,fill=memshared!20] (f1) at (7.5,3.9) {};
\node[pcell,fill=memshared!20] (f2) at (7.5,3.2) {\ldots};
\node[anchor=east] at (6.0,2.5) {k};
\node[pcell,fill=memshared!20] (fk) at (7.5,2.5) {};
\node[pcell,fill=memshared!20] (f4) at (7.5,1.8) {};
\node[pcell,fill=memshared!20] (f5) at (7.5,1.1) {};
\node[anchor=east] at (6.0,0.4) {n-1};
\node[pcell,fill=memshared!20] (fn) at (7.5,0.4) {};

% Brackets RAM / disk
\draw[decorate,decoration={brace,amplitude=4pt}]
    ($(f0.north east)+(0.15,0)$) -- ($(f2.south east)+(0.15,0)$)
    node[midway,right=5pt,font=\bfseries\small] {RAM};
\draw[decorate,decoration={brace,amplitude=4pt}]
    ($(fk.north east)+(0.15,0)$) -- ($(fn.south east)+(0.15,0)$)
    node[midway,right=5pt,font=\bfseries\small] {on disk};

\draw[-Latex,thick] (ei.east) -- ($(f1.west)+(0,0)$);
\draw[-Latex,thick] (ej.east) -- ($(fk.west)+(0,0)$);
\end{tikzpicture}
\end{center}

\begin{idea}
Cada \emph{pte} le dice al hardware ``dónde'' está realmente la página:
si en RAM (vía \texttt{ppn}) o si en disco (vía \texttt{dbn}). La RAM
se comporta entonces como una caché del conjunto total de páginas
``virtuales''. El espacio de direcciones del proceso puede ser mayor
que la RAM física.
\end{idea}

\subsection{Tratamiento del page fault}

Cuando una instrucción accede a la página $j$ con $V=0$, la CPU genera
una excepción. El kernel ejecuta un \textbf{swap-in}:

\begin{enumerate}[leftmargin=*]
    \item Reservar una página física en RAM (sea $\textit{ppn} = X$).
    \item Leer el contenido del bloque $k$ de disco en la página $X$.
    \item Mapear la \emph{pte} $j$: $\textit{ppn} = X$ y $V = 1$.
\end{enumerate}

Al retornar de la excepción, la CPU re-ejecuta la instrucción que falló
y esta vez sí puede acceder al contenido.

\subsection{Swap area}

El \textbf{swap area} es la región de disco donde se almacenan las
páginas que están fuera de RAM. Puede ser un \textbf{archivo} dentro de
un sistema de archivos común o una \textbf{partición} con un formato
específico de swap. Por ejemplo, GNU/Linux permite ambas configuraciones;
una partición con formato de tipo \emph{swap} mejora sustancialmente el
rendimiento respecto a un archivo, porque evita la indirección del
sistema de archivos.

% =====================================================================
\section{Mapeo de archivos en memoria}
% =====================================================================

Una aplicación natural y muy potente de la memoria virtual es el
\textbf{mapeo de archivos en memoria} (\emph{memory-mapped files}). La
llamada al sistema típica es:

\begin{lstlisting}[style=cstyle]
va = mmap(fd, mode);
\end{lstlisting}

\begin{definicion}
\texttt{mmap} crea un \emph{mapping} de direcciones lógicas a los
contenidos del archivo en disco y retorna la dirección lógica base del
contenido mapeado en memoria. Tras esa llamada, el proceso puede
acceder al archivo como si fuera un arreglo de bytes en memoria, sin
usar \texttt{read} o \texttt{write}.
\end{definicion}

\subsection{Cómo funciona}

El contenido del archivo se divide lógicamente en bloques del tamaño de
una página. \texttt{mmap} \emph{no} asigna memoria física ni carga el
archivo al invocarse: sólo prepara el espacio de direcciones (setea
\emph{ptes} indicando que los datos están en los bloques de disco
correspondientes, con bit $V=0$).

La carga real ocurre \textbf{bajo demanda}, de forma transparente al
proceso: cuando el programa referencia \texttt{va[i]}, el hardware
genera un page fault, el kernel asigna una página, carga en ella el
bloque correspondiente al offset $i$ y deja el \emph{pte} válido. Las
siguientes lecturas/escrituras a esa misma página ya van directo a RAM.

En general, la dirección de comienzo del mapping en sistemas tipo UNIX
es el puntero \textbf{brk} del proceso, que marca el inicio del bloque
de espacio de direcciones libre. \texttt{mmap} avanza este puntero.

Una llamada al sistema \texttt{unmmap(va)} libera el mapping y la
memoria física asociada.

\subsection{Ejemplo de uso}

\begin{lstlisting}[style=cstyle]
// myfile.dat contiene mas de 4096 bytes. Despues de mmap(),
// se ve como un arreglo de chars buf[]:
//
// buf --> +-----------------------------+ 0
//         | Hello world\0               |
//         | ...                         | 4095
//         +-----------------------------+ 4096
//         | Second page\0               | ...
//         | padding (zeroes) by mmap    |
//         +-----------------------------+ 8192

int fd = open("myfile.dat", O_RDONLY);
char *buf = mmap(fd);

printf("%s\n", buf);          // output: Hello world
printf("%s\n", buf + 4096);   // output: Second page
printf("%s\n", buf + 7005);   // output: cadena vacia (ceros)
printf("%s\n", buf + 8192);   // page fault (quiza)
\end{lstlisting}

\subsection{Ilustración del cambio en la tabla de páginas}

Antes del \texttt{mmap}, en la tabla de páginas figura el área del
proceso (text, data, stack) y, a partir de \texttt{brk}, entradas
inválidas. Después del \texttt{mmap}, las entradas que cubren la región
mapeada tienen $V=0$ con \texttt{dbn} apuntando a bloques del archivo.

\begin{center}
\begin{tikzpicture}[
    every node/.style={font=\ttfamily\small},
    sec/.style={draw,minimum width=1.9cm,minimum height=0.5cm,align=center},
    pte/.style={draw,minimum width=2.5cm,minimum height=0.5cm,align=center},
]
% Before mapping
\node[font=\bfseries\small] at (-3,3.7) {process};
\node[sec,fill=mem1!20] (t1) at (-3,3.1) {text};
\node[sec,fill=mem2!20,below=0pt of t1] (d1) {data};
\node[sec,below=0pt of d1] (f1) {free space};
\node[sec,fill=mem3!20,below=0pt of f1] (s1) {stack};
\node[anchor=east] at ($(t1.west)$) {0};
\node[anchor=east] at ($(d1.south west)+(0,-0.05)$) {brk};
\node[anchor=east] at ($(s1.south west)$) {sz};

% file
\node[draw,fill=memfile!50,minimum width=1.0cm,minimum height=0.6cm,right=0.4cm of d1] (file1) {file};

% pgtbl before
\node[font=\bfseries\small] at (-3,0.5) {page tbl};
\node[pte,fill=mem3!10] (p1a) at (-3,-0.2) {ppn0\hspace{4pt}V=1};
\node[pte,below=0pt of p1a] (p1b) {\ldots};
\node[pte,fill=memfile!10,below=0pt of p1b] (p1c) {0\hspace{4pt}V=0};
\node[pte,below=0pt of p1c] (p1d) {\ldots};
\node[pte,fill=mem3!10,below=0pt of p1d] (p1e) {ppnX\hspace{4pt}V=1};
\node[anchor=east] at ($(p1a.west)$) {0};
\node[anchor=east] at ($(p1c.west)$) {brk};
\node at (-3,-2.95) {a) Before mapping.};

% After mapping
\node[font=\bfseries\small] at (3.5,3.7) {process};
\node[sec,fill=mem1!20] (t2) at (3.5,3.1) {text};
\node[sec,fill=mem2!20,below=0pt of t2] (d2) {data};
\node[sec,fill=memfile!30,below=0pt of d2] (m2) {};
\node[sec,below=0pt of m2] (f2) {free};
\node[sec,fill=mem3!20,below=0pt of f2] (s2) {stack};
\node[anchor=east] at ($(t2.west)$) {0};
\node[anchor=east] at ($(m2.south west)$) {brk};
\node[anchor=east] at ($(s2.south west)$) {sz};

\node[draw,fill=memfile!50,minimum width=1.0cm,minimum height=0.6cm,right=0.4cm of d2] (file2) {file};
\draw[-Latex,thick] (file2) -- (m2);

% pgtbl after
\node[font=\bfseries\small] at (3.5,0.5) {page tbl};
\node[pte,fill=mem3!10] (p2a) at (3.5,-0.2) {ppn0\hspace{4pt}V=1};
\node[pte,below=0pt of p2a] (p2b) {\ldots};
\node[pte,fill=memfile!30,below=0pt of p2b] (p2c) {dbn\hspace{4pt}V=0};
\node[pte,below=0pt of p2c] (p2d) {\ldots};
\node[pte,fill=mem3!10,below=0pt of p2d] (p2e) {ppnX\hspace{4pt}V=1};
\node[anchor=east] at ($(p2a.west)$) {0};
\node[anchor=east] at ($(p2c.west)$) {va};
\node[anchor=east] at ($(p2d.west)+(0,0.15)$) {brk};
\node at (3.5,-2.95) {b) After memory mapped: va=mmap(f).};
\end{tikzpicture}
\end{center}

\subsection{Manejo del page fault sobre un mapping}

Cuando el programa accede a la región mapeada y la \emph{pte} tiene
$V=0$, el hardware genera la excepción. El kernel verifica que la
dirección lógica sea válida (por ejemplo, que el campo \texttt{ppn} no
sea cero, indicando que la entrada está \emph{configurada} para
mapping). Si así es, procesa la excepción:

\begin{enumerate}[leftmargin=*]
    \item Reservar (\emph{alloc}) una página física, sea $ppn = \alpha$.
    \item Cargar el bloque \texttt{dbn} de disco en la página obtenida.
    \item Mapear: setear $\textit{pte}.\textit{ppn} = \alpha$ y los
        flags ($V=1$) en la \emph{pte} correspondiente.
\end{enumerate}

Al retornar de la excepción, la CPU vuelve a ejecutar la instrucción que
la causó y accede exitosamente al contenido.

\paragraph{Escrituras y bit dirty.}
Si el proceso modifica el contenido en memoria del archivo, al liberar
el mapping el SO puede \emph{escribir esos datos de vuelta al archivo}.
El bit \textbf{dirty} (D) de la \emph{pte} indica qué páginas fueron
modificadas, de modo que sólo se escriben las que efectivamente
cambiaron.

% =====================================================================
\section{Carga de programas bajo demanda}
% =====================================================================

La llamada al sistema \texttt{exec(program-file, args)} debe cargar el
contenido del ejecutable a memoria. Si el programa es grande, hacerlo
de manera completa requiere tiempo y memoria, gran parte de los cuales
podrían no usarse nunca.

\begin{idea}
Si \texttt{exec} \emph{mapea} los segmentos cargables del ejecutable
con \texttt{mmap} y eventualmente carga sólo la página de código que
contiene el \emph{entry point}, es posible realizar la carga del
programa \textbf{bajo demanda}: las demás páginas se traen a RAM a
medida que la ejecución avanza y referencia nuevas instrucciones o
datos.
\end{idea}

Esta optimización es esencial en SO modernos: arrancar un binario
grande es virtualmente instantáneo porque, salvo una página, todo el
ejecutable está sólo \emph{mapeado}, no cargado.

% =====================================================================
\section{Stack con crecimiento dinámico}
% =====================================================================

Otro problema de la gestión de memoria es decidir cuánta memoria
reservar para el \emph{stack} de cada proceso. Si reservamos poco, un
programa con recursión profunda se queda sin pila; si reservamos
mucho, desperdiciamos memoria.

Con memoria virtual la solución es elegante: se mapea inicialmente
\textbf{una sola página} en una dirección lógica alta y el SO asigna y
mapea \emph{más páginas bajo demanda}, a medida que el proceso intenta
escribir sobre el tope del stack y eso genera una excepción.

\begin{center}
\begin{tikzpicture}[
    every node/.style={font=\ttfamily\small},
    sec/.style={draw,minimum width=2.0cm,minimum height=0.55cm,align=center},
]
% Antes
\node[sec,fill=mem1!20] (a1) at (0,3.5) {text};
\node[sec,fill=mem2!20,below=0pt of a1] (a2) {data};
\node[sec,fill=mem4!20,below=0pt of a2] (a3) {heap};
\node[sec,below=0pt of a3] (a4) {free};
\node[sec,below=0pt of a4] (a5) {space};
\node[sec,fill=mem3!20,below=0pt of a5] (a6) {stack};
\node[anchor=west,xshift=4pt] at (a3.east) {brk};
\node[anchor=west,xshift=4pt] at (a6.east) {sz};

% Flecha
\draw[-Latex,thick] (a6.east) ++(1.6,1.2) -- ++(1.6,0)
    node[midway,above,font=\small] {exception};

% Despues
\node[sec,fill=mem1!20] (b1) at (6.5,3.5) {text};
\node[sec,fill=mem2!20,below=0pt of b1] (b2) {data};
\node[sec,fill=mem4!20,below=0pt of b2] (b3) {heap};
\node[sec,below=0pt of b3] (b4) {free};
\node[sec,below=0pt of b4] (b5) {space};
\node[sec,fill=mem3!20,below=0pt of b5] (b6) {stack};
\node[sec,fill=mem3!35,below=0pt of b6] (b7) {stack};
\node[anchor=west,xshift=4pt] at (b3.east) {brk};
\node[anchor=west,xshift=4pt] at (b7.east) {sz};
\end{tikzpicture}
\end{center}

De esta manera el stack \emph{crece bajo demanda}. El heap,
simétricamente, suele expandirse hacia arriba (en términos lógicos)
asignándole más espacio mediante la llamada al sistema
\verb|sbrk(size)|.

% =====================================================================
\section{Memoria compartida}
% =====================================================================

Para que dos o más procesos (o el kernel y los procesos) se comuniquen
usando memoria, se requiere que \emph{uno o más bloques de memoria
físicos pertenezcan a sus espacios de direcciones simultáneamente}.
Esto se logra mapeando en cada proceso un espacio de direcciones
lógicas que resuelve a las \emph{mismas direcciones físicas}.

\begin{center}
\begin{tikzpicture}[
    every node/.style={font=\ttfamily\small},
    sec/.style={draw,minimum width=2.2cm,minimum height=0.7cm,align=center},
]
% Process 1
\node[font=\bfseries\small] at (-3,2.0) {process 1};
\node[sec] (p1a) at (-3,1.3) {\ldots};
\node[sec,fill=memshared!50,below=0pt of p1a] (p1b) {};
\node[sec,below=0pt of p1b] (p1c) {\ldots};
\node[anchor=east] at ($(p1a.west)$) {0};
\node[anchor=east] at ($(p1b.west)+(0,0.18)$) {va1};
\node[anchor=east] at ($(p1c.south west)$) {sz1};

% Shared page
\node[draw,fill=memshared!60,minimum width=2.2cm,minimum height=0.7cm] (sp) at (1.5,0.6) {shared};
\node[font=\small] at ($(sp.west)+(-0.7,0)$) {ppn=X};

\draw[-Latex,thick] (p1b.east) -- (sp.west);

% Process 2
\node[font=\bfseries\small] at (6,2.0) {process 2};
\node[sec] (p2a) at (6,1.3) {\ldots};
\node[sec,fill=memshared!50,below=0pt of p2a] (p2b) {};
\node[sec,below=0pt of p2b] (p2c) {\ldots};
\node[anchor=west] at ($(p2a.east)$) {0};
\node[anchor=west] at ($(p2b.east)+(0,0.18)$) {va2};
\node[anchor=west] at ($(p2c.south east)$) {sz2};

\draw[-Latex,thick] (p2b.west) -- (sp.east);
\end{tikzpicture}
\end{center}

Las direcciones lógicas $va_1$ y $va_2$ comúnmente son \emph{distintas}
en cada proceso, pero ambas están mapeadas al mismo número de página
física $\textit{ppn} = X$.

\subsection{Llamadas al sistema (UNIX)}

Los SO tipo UNIX proveen un conjunto de llamadas para compartir memoria:

\begin{itemize}[leftmargin=*]
    \item \verb|int shmget(key, size, flags)|: crea u obtiene un
        descriptor para un bloque de memoria compartida con
        identificador \texttt{key} de tamaño \texttt{size}.
    \item \verb|void *shmat(m)|: obtiene la dirección lógica del
        descriptor $m$ retornado por \texttt{shmget}.
    \item \verb|int shmdt(address)|: libera (\emph{detach}) el mapping.
\end{itemize}

\begin{importante}
Los procesos que se comuniquen mediante memoria compartida \emph{deben}
usar algún mecanismo de sincronización (semáforos, mutexes, \ldots) para
evitar condiciones de carrera en el acceso a los datos compartidos. La
memoria compartida \emph{no} es un canal sincronizado por sí mismo.
\end{importante}

\subsection{Bibliotecas compartidas}

Esta técnica, junto con el mapeo de archivos en memoria, es la que se
usa para mapear una \textbf{biblioteca compartida} (\emph{shared
library}) en el espacio de direcciones de un proceso. Así, varios
procesos pueden ``ver'' las mismas páginas físicas de libc, libssl,
etc., en lugar de cargar copias separadas.

\subsection{Liberación y contador de referencias}

La memoria compartida complica la liberación de memoria cuando un
proceso termina (por ejemplo, en un \texttt{exit()}): hay que liberar
sólo las páginas que \emph{no} son compartidas con otro proceso.

Una solución estándar es asociar a cada \emph{shared memory region} un
\textbf{contador de sharing} (o referencias):

\begin{itemize}[leftmargin=*,nosep]
    \item Se incrementa con cada \verb|va = shmat(m)|.
    \item Se decrementa con cada \verb|shmdt(va)|.
    \item Cuando llega a cero, es seguro liberar la memoria compartida.
\end{itemize}

% =====================================================================
\section{Copiar al escribir (Copy on Write, COW)}
% =====================================================================

La llamada al sistema \texttt{fork()} crea un proceso hijo \emph{copia}
del padre. Una implementación naïve copiaría completamente la imagen
del proceso padre en el hijo, lo cual puede requerir gran cantidad de
memoria y mucho tiempo, especialmente si luego el hijo va a hacer
\texttt{exec} y reemplazar toda esa memoria.

\subsection{La optimización}

En lugar de copiar todo, se copia sólo la \textbf{tabla de páginas},
haciendo que ambos procesos (padre e hijo) compartan inicialmente la
memoria física. Para mantener el confinamiento de cada proceso, en
\emph{ambas} tablas de páginas se quitan los permisos de
\textbf{escritura} ($W=0$) en cada \emph{pte}, dejando la memoria
compartida como \textbf{sólo de lectura}.

A cada proceso debería asociársele una \emph{shared memory block} que
describa toda su zona de direcciones COW.

\subsection{Qué pasa cuando alguien quiere escribir}

Cuando uno de los procesos intenta escribir en una página compartida
(\emph{ppn}) ocurre lo siguiente:

\begin{enumerate}[leftmargin=*]
    \item La CPU genera una excepción porque no tiene permiso de
        escritura.
    \item El kernel captura la excepción y reconoce que se debe a la
        falta de permisos, pero la dirección es \emph{válida} y
        \emph{compartida}.
    \item Asigna (\emph{alloc}) una nueva página, sea con $\textit{ppn}
        = n$.
    \item Copia el contenido de la página original en la nueva.
    \item Mapea la nueva página en el proceso que intentó escribir:
        $\textit{pte}.\textit{ppn} = n$ y $\textit{pte}.W = 1$.
    \item En el otro proceso (padre o hijo, según corresponda) también
        restaura $W=1$ en su \emph{pte}, porque ahora ese proceso es
        el único dueño de la página original.
\end{enumerate}

Como en toda excepción recuperable, el kernel prepara la CPU para que
\textbf{re-intente} la ejecución de la instrucción que provocó la
falta.

\begin{definicion}
Este mecanismo se conoce como \textbf{copy on write} (\textbf{COW}): la
copia real se difiere hasta el momento en que efectivamente se necesita
porque alguien escribe. Mientras los procesos sólo lean, no hay
duplicación de memoria.
\end{definicion}

Para generalizar COW se necesita llevar la pista de las páginas
compartidas entre procesos, comúnmente usando \emph{shared memory
regions} en los descriptores de tareas (los \emph{task structs} del
kernel).

\begin{idea}
COW es lo que hace que \texttt{fork()} + \texttt{exec()} sea barato: el
padre no se ``clona'' físicamente; sólo se prepara para clonarse si
hace falta. Como \texttt{exec()} suele descartar toda la memoria del
hijo, la copia nunca llega a producirse.
\end{idea}

% =====================================================================
\section{Algoritmos de reemplazo de páginas}
% =====================================================================

Cuando el sistema requiere asignar memoria (por ejemplo
\verb|allocpage()|), ya sea para un proceso o para el kernel, puede
ocurrir que no exista ninguna página libre. Esta situación se conoce
como \textbf{page fault}.

Con swapping el sistema puede elegir una página (la \textbf{víctima})
que actualmente está en uso, hacerle un \emph{swap-out} a disco y
reutilizar el frame para satisfacer el requerimiento.

\begin{importante}
El problema central es: \textbf{¿cómo elegir la víctima?} Lo ideal
sería que fuera una página que \emph{no} va a referenciarse en mucho
tiempo (o nunca más). Pero esto no es computable: requiere predecir el
futuro. Los algoritmos prácticos intentan aproximarlo con heurísticas
basadas en el comportamiento pasado.
\end{importante}

% ---------------------------------------------------------------------
\subsection{FIFO}
% ---------------------------------------------------------------------

El sistema mantiene una cola de las páginas en uso: la última asignada
queda al final. Ante un page fault, el algoritmo \textbf{FIFO} elige
como víctima la página a la \emph{cabeza} de la cola (la que entró hace
más tiempo). La nueva página se encola al final.

\paragraph{Crítica.}
FIFO es muy simple, pero puede elegir como víctima una página que se
usa con \textbf{altísima frecuencia} (por ejemplo, la página de código
con el bucle principal del programa), sólo por haber entrado primero.

% ---------------------------------------------------------------------
\subsection{Second chance (algoritmo del reloj, \emph{clock})}
% ---------------------------------------------------------------------

Es una corrección sencilla de FIFO que aprovecha el bit \textbf{A}
(\emph{accessed}) que el hardware setea cuando una página es leída o
escrita.

En vez de tomar la página de la cabeza directamente, se inspecciona
primero su bit \emph{A}:

\begin{itemize}[leftmargin=*]
    \item Si $A=1$, se hace $A=0$ y la página se mueve al final de la
        cola, dándole una ``segunda oportunidad''.
    \item Si $A=0$, se la elige como víctima.
\end{itemize}

Esto continúa hasta encontrar una página a la cabeza con $A=0$. La idea
es que la víctima elegida \emph{no} fue accedida desde el último
recorrido.

\paragraph{Variante: el algoritmo del reloj.}
En lugar de usar una cola, se usa una \textbf{lista circular} de las
páginas usadas, con un puntero o \textbf{hand} que apunta a la página
``más vieja'' candidata. Ante un page fault se verifica el bit $A$ de
la página apuntada por \emph{hand}: si es cero, se la elige; si no, se
hace $A=0$ y se avanza \emph{hand} repitiendo el proceso. Esta
implementación se conoce como \textbf{algoritmo del reloj}
(\emph{clock}).

% ---------------------------------------------------------------------
\subsection{Menos recientemente usada (LRU)}
% ---------------------------------------------------------------------

\textbf{LRU} (\emph{Least Recently Used}) busca como víctima una página
que no ha sido usada en el último lapso de tiempo.

\paragraph{Idea pura (costosa).}
Mantener una lista ordenada de todas las páginas por tiempo de acceso
(de menor a mayor). En cada intervalo de tiempo (\emph{ticks}), las
páginas accedidas se mueven al final. Ante un page fault, la víctima es
la cabeza. Sin soporte de hardware específico, esta implementación es
inviable: cada acceso a memoria debería actualizar la lista.

\paragraph{Implementación práctica: NFU.}
Una implementación alternativa se conoce como \textbf{Not Frequently
Used} (\textbf{NFU}):

\begin{itemize}[leftmargin=*]
    \item A cada página se le asocia un \emph{contador} inicialmente
        en cero.
    \item Periódicamente, el SO recorre cada \emph{pte} de cada tabla
        de páginas y \emph{suma el bit $A$ al contador}.
    \item Se mantiene una lista ordenada por contador, de menor a
        mayor.
    \item Ante un page fault, la víctima es la página a la cabeza
        (menor contador).
\end{itemize}

\paragraph{Dos problemas de NFU y sus soluciones.}

\begin{enumerate}[leftmargin=*]
    \item \textbf{Recorrer todas las tablas periódicamente es caro.}
        Se mitiga haciéndolo en momentos de \emph{ocio} del sistema y
        en forma \emph{incremental}.
    \item \textbf{El contador sólo crece: nunca olvida.} Páginas que
        fueron muy usadas en el pasado pero ahora no lo son, quedarían
        protegidas para siempre. Se resuelve con \textbf{aging}: antes
        de sumarle el bit $A$, se hace un \emph{shift a derecha} del
        contador (dividir por dos). Así, los accesos antiguos pesan
        cada vez menos.
\end{enumerate}

\paragraph{Caso de estudio: Linux.}
El kernel de Linux usa actualmente una versión de \emph{software LRU}
(NFU) de \textbf{dos niveles}, manteniendo dos listas:
\textbf{active\_list} (páginas accedidas recientemente) e
\textbf{inactive\_list}. Las páginas inactivas son las candidatas a ser
reemplazadas.

% ---------------------------------------------------------------------
\subsection{No recientemente usada (NRU)}
% ---------------------------------------------------------------------

\textbf{NRU} (\emph{Not Recently Used}) selecciona una página a
reemplazar usando exclusivamente el soporte de los bits $A$ (accessed)
y $D$ (dirty) del hardware.

El objetivo es identificar qué páginas fueron accedidas en el último
período de tiempo. Periódicamente, cada $n$ ticks, el bit $A$ de cada
\emph{pte} se pone en cero (\emph{cleared}).

\paragraph{Clasificación en clases.}
Ante un page fault, las páginas se clasifican en 4 clases:

\begin{enumerate}[leftmargin=*]
    \item $A=0,\;D=0$: no accedida en el último período (mejor
        candidata).
    \item $A=0,\;D=1$: escrita en el período anterior; el $A=0$ es
        producto del \emph{clearing} periódico.
    \item $A=1,\;D=0$: leída en el último período.
    \item $A=1,\;D=1$: accedida en el último período, y además fue
        escrita en algún momento (peor candidata).
\end{enumerate}

Se elige una víctima en la \textbf{menor clase no vacía}. Este algoritmo
se conoce también como \emph{software LRU} porque permite implementar
una aproximación de LRU en arquitecturas modernas \emph{sin} soporte
específico para LRU.

% ---------------------------------------------------------------------
\subsection{Working set}
% ---------------------------------------------------------------------

\begin{definicion}
El \textbf{working set} de un proceso es el conjunto de páginas a las
que el proceso accedió en un \emph{período de tiempo} reciente $\tau$.
\end{definicion}

\begin{idea}
Los programas raramente acceden a su memoria al azar: en una iteración
o cuando trabajan sobre datos contiguos, referencian \emph{muchas
veces} un mismo conjunto de páginas. Mantener ese conjunto en RAM y
desalojar el resto es una excelente heurística.
\end{idea}

\paragraph{Implementación.}
Se asocia a cada página un \emph{reloj lógico} (en ticks) que registra
el tiempo del último acceso (\emph{tick}).

Ante un page fault se recorre la tabla de páginas inspeccionando el
bit $A$ de cada \emph{pte}:

\begin{itemize}[leftmargin=*]
    \item Si $A=1$, se actualiza el reloj asociado a la página y se
        agrega al working set del período corriente.
    \item Si $A=0$ se evalúa si puede elegirse como víctima:
        \begin{itemize}[leftmargin=*,nosep]
            \item Si su tiempo de acceso cae dentro del intervalo
                actual ($\textit{tick} - \textit{last\_access}_{\textit{page}} < \tau$),
                se mantiene en el working set.
            \item Si no, se considera \emph{fuera} del working set.
        \end{itemize}
\end{itemize}

Se elige como víctima una página fuera del working set con menor valor
de reloj.

Para evitar recorrer toda la tabla de páginas en cada falta, puede
combinarse con una mejora del estilo del \emph{algoritmo del reloj}
descripto antes.

El algoritmo de Linux, descripto en \S\,\ref{sec:linux-mv}, intenta de
alguna manera que la lista de páginas \emph{activas} represente el
working set actual del sistema.

% ---------------------------------------------------------------------
\subsection{Otras consideraciones}
% ---------------------------------------------------------------------

Los algoritmos descriptos arriba no consideran ciertos aspectos
adicionales que el SO real sí debe tener en cuenta:

\paragraph{Política global vs.\ local.}
Hay que decidir si la política de reemplazo es:

\begin{itemize}[leftmargin=*,nosep]
    \item \textbf{Global}: no se considera a qué proceso pertenece la
        página candidata.
    \item \textbf{Local}: sólo se buscan víctimas entre las páginas
        \emph{del proceso corriente}.
\end{itemize}

No es lo mismo reemplazar una página del kernel o de una biblioteca
compartida, que una página de un proceso de usuario.

\paragraph{Páginas protegidas.}
Para evitar problemas, ciertas páginas se \textbf{protegen} y se marcan
para que no puedan ser elegidas como víctimas. Por ejemplo, las páginas
del kernel o las que pertenecen a espacios de direcciones de dispositivos
mapeados en memoria (\textbf{MMIO}), donde escribir/leer tiene efectos
físicos sobre el hardware.

\paragraph{Trashing (sobrepaginación).}
Uno de los peligros del \emph{page swapping} es justamente este
fenómeno.

\begin{importante}
La \textbf{sobrepaginación} (\emph{trashing}) es un estado en el que el
sistema está corriendo \emph{out of memory} y ocurren repetidamente
los siguientes eventos:
\begin{enumerate}[leftmargin=*,nosep]
    \item Se hace un swap-out para liberar memoria.
    \item La página recién desalojada vuelve a ser referenciada y hay
        que cargarla de nuevo (swap-in).
\end{enumerate}
\end{importante}

Estos pasos se repiten a alta frecuencia; el sistema degrada el
rendimiento drásticamente y el progreso de los procesos se hace muy
lento.

% =====================================================================
\section{Caso de estudio: GNU/Linux}
\label{sec:linux-mv}
% =====================================================================

Linux implementa un algoritmo \textbf{LRU}, aunque la lista de páginas
no refleja exactamente un orden estricto por tiempo de acceso.

Mantiene dos listas:

\begin{itemize}[leftmargin=*]
    \item \texttt{active\_list}: idealmente contiene un \emph{working
        set} de todos los procesos.
    \item \texttt{inactive\_list}: lista de candidatos a ser
        reemplazados.
\end{itemize}

Estas listas se actualizan periódicamente. Esto hace que la política
sea \textbf{global}: cualquier página candidata puede ser elegida,
independientemente de a qué proceso pertenece.

% =====================================================================
\section{Cierre}
% =====================================================================

La memoria virtual es la pieza que conecta el paginado del hardware con
la realidad de los sistemas modernos, donde la RAM siempre es menos
abundante que la memoria que los procesos pretenden usar.

\begin{itemize}[leftmargin=*]
    \item El \textbf{swapping / demand paging} permite que la RAM se
        comporte como una caché del espacio de direcciones lógicas: el
        bit $V$ y el \emph{disk block number} en la \emph{pte} indican
        si una página vive en RAM o en disco.
    \item El \textbf{mapeo de archivos} (\texttt{mmap}) usa la misma
        maquinaria para exponer archivos como arreglos de memoria. La
        carga se difiere hasta el primer acceso, y el bit \emph{dirty}
        marca qué páginas hay que escribir de vuelta al cerrar.
    \item La \textbf{carga bajo demanda} y el \textbf{crecimiento
        dinámico del stack} reducen drásticamente el costo de
        \texttt{exec()} y eliminan la necesidad de reservar pila por
        adelantado.
    \item La \textbf{memoria compartida} (\texttt{shmget}, \texttt{shmat},
        \texttt{shmdt}) y las bibliotecas compartidas se apoyan en
        mapear el mismo \texttt{ppn} en varias tablas de páginas; un
        contador de referencias resuelve la liberación segura.
    \item El \textbf{copy on write} es lo que hace que \texttt{fork()}
        sea barato en la práctica: el padre y el hijo comparten la
        memoria en sólo-lectura hasta que alguien escribe.
    \item Los \textbf{algoritmos de reemplazo} (FIFO, second
        chance/clock, LRU/NFU con aging, NRU por clases, working set)
        son aproximaciones distintas al mismo problema imposible:
        adivinar qué página no se va a usar más. Cada uno equilibra la
        precisión de la heurística contra el costo de implementarla.
    \item \textbf{Linux} combina varias de estas ideas en un esquema de
        dos listas (\texttt{active}/\texttt{inactive}) con política
        global, intentando aproximar el working set del sistema.
\end{itemize}

\begin{idea}
La memoria virtual no es ``un truco'' aislado: es una capa que se
construye sobre el paginado y vuelve a usar las mismas herramientas
(bits $V$, $A$, $D$, \emph{page faults}) para resolver problemas muy
distintos: ejecutar procesos grandes, compartir bibliotecas, abaratar
\texttt{fork}, mapear archivos, hacer crecer el stack. El próximo
capítulo (\emph{Entrada/Salida}) se apoya también en muchas de estas
ideas, especialmente en el mapeo en memoria de dispositivos.
\end{idea}

\end{document}
